\documentclass{beamer}

\usepackage[utf8]{inputenc}
\usepackage[T5]{fontenc}
\usepackage[vietnamese]{babel}

\usepackage{amsmath,amssymb}
\usepackage{tikz}
\usepackage{listings}
\usepackage{array}
\usepackage{colortbl}
\usepackage{booktabs}
\usepackage{xcolor}
\usepackage{tcolorbox}

\usetheme{Madrid}
\usecolortheme{dolphin}

% --- màu ---
\definecolor{myblue}{RGB}{30,144,255}
\definecolor{mygreen}{RGB}{60,179,113}
\definecolor{myred}{RGB}{220,20,60}
\definecolor{mypurple}{RGB}{138,43,226}
\definecolor{myorange}{RGB}{255,140,0}
\definecolor{mycyan}{RGB}{0,206,209}
\definecolor{codebg}{RGB}{248,250,252}
\definecolor{codeid}{RGB}{255,140,0}

\setbeamercolor{structure}{fg=myblue}
\setbeamercolor{frametitle}{fg=white,bg=myblue}
\setbeamerfont{frametitle}{series=\bfseries}

% --- box đẹp ---
\tcbset{
colback=blue!5,
colframe=myblue,
arc=3mm,
boxrule=0.8pt
}

% macros
\newcommand{\hlblue}[1]{\textcolor{myblue}{#1}}
\newcommand{\hlgreen}[1]{\textcolor{mygreen}{#1}}
\newcommand{\hlred}[1]{\textcolor{myred}{#1}}
\newcommand{\hlpurple}[1]{\textcolor{mypurple}{#1}}
\newcommand{\hlorange}[1]{\textcolor{myorange}{#1}}
\newcommand{\hlcyan}[1]{\textcolor{mycyan}{#1}}
\newcommand{\hlwhite}[1]{\textcolor{white}{#1}}

% --- code đẹp ---
\lstset{
basicstyle=\ttfamily\small,
backgroundcolor=\color{codebg},
keywordstyle=\color{myblue}\bfseries,
commentstyle=\color{mygreen}\itshape,
stringstyle=\color{mypurple},
identifierstyle=\color{codeid},
numbers=left,
numberstyle=\tiny\color{gray},
breaklines=true,
showstringspaces=false,
frame=single,
rulecolor=\color{myblue}
}

% --- title đẹp ---
\setbeamertemplate{title page}
{
\centering
\vspace{2cm}
{\usebeamerfont{title}\color{myred}\inserttitle\par}
\vspace{1cm}
{\usebeamerfont{author}\insertauthor\par}
\vspace{0.5cm}
{\usebeamerfont{date}\insertdate\par}
}

\title{\hlred{DAG - Directed Acyclic Graph}}
\author{Group 3}
\date{\hlcyan{03/2026}}

\begin{document}

\frame{\titlepage}

% ===================== LÝ THUYẾT =====================

\begin{frame}{\hlwhite{Định nghĩa DAG}}
\begin{tcolorbox}[title=\textbf{\hlwhite{Khái niệm DAG}}]
\hlblue{DAG} là đồ thị:
\begin{itemize}
\item Có hướng
\item Không chứa chu trình
\end{itemize}
\end{tcolorbox}

\begin{tcolorbox}
\centering
\hlpurple{$v \rightarrow ... \rightarrow v$ không tồn tại}
\end{tcolorbox}
\end{frame}

\begin{frame}{\hlwhite{Bản chất của DAG}}
\begin{itemize}
\item Không có chu trình $\rightarrow$ không quay lại điểm ban đầu
\item Quan hệ giữa các đỉnh:
\begin{itemize}
\item \hlblue{Thứ tự}
\item \hlblue{Phụ thuộc}
\end{itemize}
\end{itemize}

\begin{tcolorbox}
\centering
\hlpurple{DAG $\Rightarrow$ có thứ tự tuyến tính}
\end{tcolorbox}
\end{frame}

\begin{frame}{\hlwhite{Định lý quan trọng}}
\begin{tcolorbox}[title=\textbf{\hlwhite{Định lý}}]
Đồ thị là DAG \hlred{khi và chỉ khi} có thể topo sort
\end{tcolorbox}

\begin{itemize}
\item Topo thành công $\Rightarrow$ DAG
\item Topo Thất bại $\Rightarrow$ có chu trình
\end{itemize}
\end{frame}

\begin{frame}{\hlwhite{Tính chất}}
\begin{itemize}
\item Luôn tồn tại đỉnh:
\begin{itemize}
\item số cạnh đi vào đỉnh đó: \hlblue{in-degree = 0} (điểm đầu)
\item số cạnh đi ra từ đỉnh đó: \hlgreen{out-degree = 0} (điểm cuối)
\end{itemize}
\item Không có chu trình
\item Không lặp vô hạn
\end{itemize}
\end{frame}

% ===================== MINH HỌA =====================

\begin{frame}{\hlwhite{Ví dụ DAG}}
\centering
\begin{tikzpicture}
\tikzset{v/.style={circle,draw=myblue,fill=myblue!10,minimum size=8mm}}

\node[v] (1) at (0,0) {1};
\node[v] (2) at (2,1.5) {2};
\node[v] (3) at (4,0) {3};
\node[v] (4) at (6,1.5) {4};

\draw[->] (1)--(2);
\draw[->] (1)--(3);
\draw[->] (2)--(4);
\draw[->] (3)--(4);
\end{tikzpicture}
\end{frame}

\begin{frame}{\hlwhite{Không phải DAG}}
\centering
\begin{tikzpicture}
\tikzset{v/.style={circle,draw=myred,fill=myred!10,minimum size=8mm}}

\node[v] (1) at (0,0) {1};
\node[v] (2) at (2,1.5) {2};
\node[v] (3) at (4,0) {3};

\draw[->] (1)--(2);
\draw[->] (2)--(3);
\draw[->] (3)--(1);
\end{tikzpicture}
\end{frame}

\end{document}